\chapter{Dissociation}

\section*{Definition}
Dissociation is a disruption in the normal integration of consciousness, memory, identity, emotion, perception, body representation, motor control, or behavior. Dissociative disorders include depersonalization/derealization disorder, dissociative amnesia, and dissociative identity disorder.

\section*{What Happens in Your Body}
Dissociation reflects a failure of integrative brain function, often in the context of overwhelming stress or trauma. The prefrontal cortex exerts excessive inhibitory control over limbic structures (amygdala, insula), resulting in emotional suppression and detachment, while hippocampal dysfunction impairs memory integration and encoding.

\section*{Causes}
\begin{itemize}
  \item \textbf{Psychiatric:} Dissociative disorders, PTSD, acute stress disorder, borderline personality disorder, conversion disorder
  \item \textbf{Trauma:} Childhood physical, sexual, or emotional abuse; neglect; combat exposure; natural disasters; torture
  \item \textbf{Neurological:} Temporal lobe epilepsy, traumatic brain injury, migraine, brain tumor
  \item \textbf{Substance-induced:} Ketamine, cannabis, hallucinogens, alcohol
  \item \textbf{Cultural:} Trance and possession states (culturally sanctioned dissociative phenomena)
\end{itemize}

\section*{When to See a Doctor}
See your doctor if:
\begin{itemize}
  \item Dissociative episodes are recurrent, prolonged, or cause distress and functional impairment
  \item There are memory gaps for extended periods (hours to days)
  \item The individual finds evidence of actions they do not recall (suggesting dissociative identity disorder)
\end{itemize}

\section*{Related Symptoms}
\begin{itemize}
  \item Depersonalization, derealization, dissociative amnesia, identity confusion, trance states, conversion symptoms, PTSD, self-harm, suicidal ideation
\end{itemize}
\textbf{Important:} Dissociation exists on a continuum from normal absorption (daydreaming) to pathological loss of memory and identity.

\section*{Self-Care / Home Management}
\begin{itemize}
  \item Use grounding techniques: temperature sensations, object descriptions, physical movement
  \item Keep a symptom diary to identify triggers and dissociative episodes
  \item Establish a regular routine with consistent sleep, meals, and activity
  \item Avoid alcohol, cannabis, and stimulants that may worsen dissociation
\end{itemize}

\section*{How Your Doctor Will Evaluate You}
\begin{itemize}
  \item Clinical interview with DSM-5 diagnostic criteria
  \item Structured instruments: Dissociative Experiences Scale (DES), Structured Clinical Interview for DSM-5 Dissociative Disorders (SCID-D)
  \item Rule out medical causes: EEG to exclude temporal lobe epilepsy; neuroimaging if focal signs
  \item Screen thoroughly for trauma history, PTSD, and suicidal ideation
\end{itemize}

\section*{Treatment Options}
\begin{itemize}
  \item Phase-oriented psychotherapy: (1) stabilization and safety, (2) trauma processing, (3) identity integration and rehabilitation
  \item Trauma-focused therapy: EMDR, cognitive processing therapy, prolonged exposure
  \item No FDA-approved pharmacotherapy; treat comorbid depression, anxiety, or PTSD with SSRIs
  \item Hospitalization if self-harm, suicidality, or severe functional impairment
\end{itemize}

\section*{References}
\begin{thebibliography}{99}
\bibitem{ref1} American Psychiatric Association. "Dissociative disorders." In: Diagnostic and Statistical Manual of Mental Disorders. 5th ed, Text Revision. APA; 2022.
\bibitem{ref2} Spiegel D, Loewenstein RJ, Lewis-Fernandez R, et al. "Dissociative disorders in DSM-5." Depress Anxiety. 2011;28(12):E17-E29.
\bibitem{ref3} Brand BL, Lanius RA, Vermetten E, et al. "Dissociative disorders: an update." Harv Rev Psychiatry. 2014;22(4):225-238.
\bibitem{ref4} Sar V. "Dissociative disorders: an overview." J Trauma Dissociation. 2010;11(3):247-262.
\bibitem{ref5} Howell EF. "The Dissociative Mind." Routledge; 2011.
\bibitem{ref6} Spiegel D. "Dissociative disorders: clinical features and diagnosis." UpToDate. Wolters Kluwer; 2025.
\end{thebibliography}

\clearpage
